% \ljm{where's the binding constraint here?}

% Let $\mathcal{M}$ be the space of language models and $\mathcal{D}$ the set of deployment contexts.
% Each context $d \in \mathcal{D}$ is characterized by:

% \begin{itemize}[leftmargin=5mm,topsep=0mm,itemsep=0mm]
%   \item Resource availability, $\mathbf{r}_d \in \mathbb{R}_{\geq 0}^k$, representing the resources available for deployment.
%   \item Performance function, $P_d: \mathcal{M} \rightarrow \mathbb{R}$, capturing task-aligned utility.
%   \item Cost function, $\mathbf{c}_d: \mathcal{M} \rightarrow \mathbb{R}_{\geq 0}^k$, capturing deployment requirements.
% \end{itemize}

% For a given context $d$, the best achievable performance is:

% \vspace{-1em}
% \begin{equation}
%   P_d^{\star} \max_{m \in \mathcal{M}} P_d(m)~\text{s.t.} ~\mathbf{c}_d(m) \leq \mathbf{r}_d
% \end{equation}
% \vspace{-1em}

% The Global Access Frontier is the mapping from resource availability to optimal performance across deployment contexts:

% \vspace{-1em}
% \begin{equation}
%   \mathcal{F}: \mathbf{r}_d \mapsto P^{\star}_d
% \end{equation}
% \vspace{-1em}

% or equivalently, the set of solutions $\mathcal{F} = \{(\mathbf{r}_d, P^{\star}_d): d \in \mathcal{D}\}$.

% \paragraph{Efficiency Frontier is a special case}
% The Efficiency Frontier assumes a single dimension of constraint, i.e., compute.
% We also argue that this frontier assumes context-independent performance and cost:

% \vspace{-1em}
% \begin{equation}
%   P^{\star} = \max_{m \in \mathcal{M}} P(m) \text{s.t.} c(m) \leq r
% \end{equation}
% \vspace{-1em}

% This is equivalent to setting $k=1$, $P_d=P~\forall~d$, and $\mathbf{c}_d=c~\forall~d$.
% Although this formulation has driven significant progress, it leaves a question unasked: \textit{efficient for whom?}
% A model that achieves state-of-the-art performance-per-FLOP may still be inaccessible to communities constrained by unreliable infrastructure or the absence of meaningful evaluation.
% The Global Access Frontier generalizes by treating all constraints explicitly.

% \subsection{Deployment Contexts}

% We describe a deployment context $d \in \mathcal{D}$

% % we now discuss what these deployment contexts are
% % add a table

% \paragraph{The ``goldilocks'' zone}
% % demand-side changes

% \paragraph{Recent progress in NLP}

% \subsection{Binding Constraints}

% \paragraph{Access to Data}


% \paragraph{Access to Compute}


% \paragraph{Access to Benchmarks}


% \paragraph{Access to Infrastructure}


% \paragraph{Access to Expertise}

\subsection{Problem Formulation and Notation}

We formalize the Global Access Frontier by drawing on the theory of binding constraints from growth diagnostics \citep{hausmann2005growth}. The key insight is that deployment contexts face \emph{multiple} constraints simultaneously, but not all constraints bind equally---identifying and relaxing the most binding constraint yields the largest gains.

\paragraph{Setup.} Let $\mathcal{M}$ be the space of language models and $\mathcal{D}$ the set of deployment contexts. Each context $d \in \mathcal{D}$ is characterized by:
\begin{itemize}
    \item \textbf{Resource availability} $r_d = (r_d^{(1)}, \ldots, r_d^{(k)}) \in \mathbb{R}^k_{\geq 0}$, representing available resources across $k$ dimensions (e.g., compute, data, infrastructure, benchmarks, expertise).
    \item \textbf{Performance function} $P_d: \mathcal{M} \to \mathbb{R}$, capturing task-aligned utility for context $d$.
    \item \textbf{Cost functions} $c_d^{(i)}: \mathcal{M} \to \mathbb{R}_{\geq 0}$ for $i \in \{1, \ldots, k\}$, where $c_d^{(i)}(m)$ is the consumption of resource $i$ when deploying model $m$ in context $d$.
\end{itemize}

\paragraph{Constrained Optimization.} For a given context $d$, the best achievable performance is:
\begin{equation}
    P_d^\star = \max_{m \in \mathcal{M}} P_d(m) \quad \text{s.t.} \quad c_d^{(i)}(m) \leq r_d^{(i)} \quad \forall i \in \{1, \ldots, k\}
    \label{eq:constrained-opt}
\end{equation}

\paragraph{Binding Constraints and Shadow Prices.} Let $\lambda_d^{(i)} \geq 0$ denote the Lagrange multiplier (shadow price) associated with constraint $i$ in context $d$. By the envelope theorem:
\begin{equation}
    \frac{\partial P_d^\star}{\partial r_d^{(i)}} = \lambda_d^{(i)}
    \label{eq:shadow-price}
\end{equation}
The shadow price $\lambda_d^{(i)}$ measures the marginal value of relaxing constraint $i$: how much task-aligned performance improves per unit of additional resource $i$.

\paragraph{Definition 1 (Binding Constraint).} Constraint $i$ is \emph{binding} in context $d$ if: (1) the constraint is active, i.e., $c_d^{(i)}(m^\star) = r_d^{(i)}$, and (2) the shadow price is positive, i.e., $\lambda_d^{(i)} > 0$. Constraint $i$ is the \emph{most binding constraint} if $\lambda_d^{(i)} = \max_{j} \lambda_d^{(j)}$.

The diagnostic question central to the Global Access Frontier is: \emph{for a given deployment context $d$, which constraint is most binding?} Formally:
\begin{equation}
    i^\star_d = \argmax_{i \in \{1, \ldots, k\}} \lambda_d^{(i)}
    \label{eq:diagnostic}
\end{equation}

\paragraph{The Global Access Frontier.} We define the Global Access Frontier as the mapping from resource profiles to optimal performance across all deployment contexts:
\begin{equation}
    \mathcal{F} = \left\{ \left( r_d, P_d^\star, i_d^\star \right) : d \in \mathcal{D} \right\}
    \label{eq:frontier}
\end{equation}
This characterizes not only what performance is achievable (as in the Efficiency Frontier), but also \emph{why} certain contexts fall short---by identifying their binding constraints.

\paragraph{The Efficiency Frontier as a Special Case.} The Efficiency Frontier assumes a single constraint dimension (compute) and context-independent performance:
\begin{equation}
    P^\star = \max_{m \in \mathcal{M}} P(m) \quad \text{s.t.} \quad c(m) \leq r
    \label{eq:efficiency-frontier}
\end{equation}
This is equivalent to setting $k = 1$, $P_d = P$ for all $d$, and $c_d = c$ for all $d$. While this formulation has driven significant progress in efficient NLP, it leaves a critical question unasked: \emph{efficient for whom?} A model achieving state-of-the-art performance-per-FLOP may remain inaccessible to communities constrained by data scarcity or infrastructure limitations. The Global Access Frontier generalizes by treating all constraints explicitly and identifying which one binds.

% ==============================================================================
\subsection{Deployment Contexts}
% ==============================================================================

A deployment context $d \in \mathcal{D}$ represents a specific setting in which an LLM may be deployed, characterized by its resource profile $r_d$ and performance requirements. We identify five primary constraint dimensions:

\begin{enumerate}
    \item \textbf{Compute} ($r_d^{\text{compute}}$): Available computational resources, including hardware capabilities, memory, and energy budget.
    \item \textbf{Data} ($r_d^{\text{data}}$): Availability of training and adaptation data in the target language(s) and domain(s).
    \item \textbf{Infrastructure} ($r_d^{\text{infra}}$): Network connectivity, latency tolerance, and deployment platform constraints.
    \item \textbf{Benchmarks} ($r_d^{\text{bench}}$): Availability of meaningful evaluation metrics and test sets for the target task and language.
    \item \textbf{Expertise} ($r_d^{\text{expert}}$): Local technical capacity to deploy, adapt, and maintain models.
\end{enumerate}

Different deployment contexts exhibit markedly different resource profiles. Table~\ref{tab:contexts} illustrates this heterogeneity.

\begin{table}[t]
    \centering
    \small
    \begin{tabular}{lccccc}
        \toprule
        \textbf{Context} & \textbf{Compute} & \textbf{Data} & \textbf{Infra} & \textbf{Bench} & \textbf{Expert} \\
        \midrule
        Hyperscaler      & High             & High          & High           & High           & High            \\
        Enterprise       & High             & Med           & High           & Med            & Med             \\
        Industry Lab     & Med              & Med           & High           & Med            & High            \\
        University       & Med              & Low           & Med            & Low            & High            \\
        Individual       & Low              & Low           & Med            & Low            & Low             \\
        Rural Community  & Low              & Low           & Low            & Low            & Low             \\
        \bottomrule
    \end{tabular}
    \caption{Illustrative resource profiles across deployment contexts. The binding constraint varies: for hyperscalers, no constraint strongly binds; for rural communities, infrastructure and data are likely most binding.}
    \label{tab:contexts}
\end{table}

% ==============================================================================
\subsection{Binding Constraints}
% ==============================================================================

Following \citet{hausmann2005growth}, we operationalize binding constraints through observable symptoms. A constraint is likely binding if:

\begin{enumerate}
    \item \textbf{High shadow price}: Agents expend significant resources to circumvent the constraint (e.g., communities building custom datasets signal data scarcity).
    \item \textbf{Responsiveness}: Relaxing the constraint produces measurable performance gains (e.g., adding target-language data improves downstream accuracy).
    \item \textbf{Bypassing behavior}: Economic actors find workarounds (e.g., code-switching to English when native-language models fail).
    \item \textbf{Differential thriving}: Entities less reliant on the constrained resource outperform those that depend on it.
\end{enumerate}

Table~\ref{tab:diagnostics} maps observable symptoms to likely binding constraints in the LLM deployment setting.

\begin{table}[t]
    \centering
    \small
    \begin{tabular}{p{5.5cm}l}
        \toprule
        \textbf{Symptom}                                          & \textbf{Likely Binding Constraint} \\
        \midrule
        High English performance, low target-language performance & Data                               \\
        Model exists but cannot run locally                       & Compute / Infrastructure           \\
        No reliable way to evaluate task performance              & Benchmarks                         \\
        Model deployed but users cannot use it effectively        & Expertise                          \\
        Significant latency or connectivity failures              & Infrastructure                     \\
        Communities building own datasets                         & Data                               \\
        Reliance on English as intermediary                       & Data / Benchmarks                  \\
        \bottomrule
    \end{tabular}
    \caption{Diagnostic mapping: observable symptoms and their likely binding constraints.}
    \label{tab:diagnostics}
\end{table}

% ==============================================================================
\subsection{The Goldilocks Zone}
% ==============================================================================

We define the \emph{Goldilocks zone} as the region where a model is both deployable and useful for a given context. Let $\tau_d$ denote the minimum task-aligned performance required for meaningful utility in context $d$. Then:

\begin{equation}
    \mathcal{G} = \left\{ (m, d) \in \mathcal{M} \times \mathcal{D} : P_d(m) \geq \tau_d \;\land\; c_d^{(i)}(m) \leq r_d^{(i)} \;\forall i \right\}
    \label{eq:goldilocks}
\end{equation}

The Goldilocks zone captures deployment configurations where progress is \emph{immediately useful}: models that meet both the performance threshold and resource constraints of a given context.

\paragraph{Expanding the Frontier.} Research progress expands the Goldilocks zone in different directions (Figure~\ref{fig:frontier}):
\begin{itemize}
    \item \textbf{Efficient NLP} (quantization, distillation, pruning): Reduces $c_d^{\text{compute}}(m)$, expanding access for compute-constrained contexts.
    \item \textbf{Multilingual NLP} (cross-lingual transfer, multilingual pretraining): Increases $P_d(m)$ for underserved languages, expanding access for data-constrained contexts.
    \item \textbf{Infrastructure-aware deployment}: Reduces $c_d^{\text{infra}}(m)$ through offline capabilities, edge optimization, etc.
\end{itemize}

The Global Access Frontier framework clarifies that these research directions are \emph{complementary}: each relaxes a different constraint, and the relevant direction depends on which constraint binds for a given context.


\begin{comment}
1. Introduction

Current efficiency frontier optimizes capability per FLOP
But this is context-blind: doesn't ask who benefits
We propose a "global access frontier" that centers deployment constraints and usefulness
Preview the taxonomy

2. Background: Growth Diagnostics for LLM Access

Development economics rejected one-size-fits-all prescriptions (Washington Consensus)
Hausmann, Rodrik, & Velasco (2005): identify binding constraints that differ by context
We apply this logic to LLM deployment: different contexts face different bottlenecks
The taxonomy organizes methods by which constraint they relax

3. Supply-side: Model Efficiency

Architecture efficiency (efficient attention, SSMs, MoE)
Compression (quantization, pruning, distillation)
Adaptive computation (early exit, speculative decoding)

4. Supply-side: Data Efficiency

Low-resource language modeling
Cross-lingual transfer
Parameter-efficient fine-tuning (LoRA, adapters)

5. Demand-side: Task Alignment

Mismatch between benchmarks and actual user needs
Localized evaluation
Domain adaptation for underserved contexts

6. Infrastructure: Deployment Systems

Edge runtimes, on-device inference
Offline-first architectures
Hybrid cloud-edge systems

7. Institutional: Access Mechanisms

Open weights
Pricing and API economics
Licensing and governance

8. Discussion / Open Problems

Gaps in current research
Misaligned incentives
Toward access-aware evaluation

9. Conclusion
\end{comment}

